\documentclass[11pt,a4paper]{book}
\usepackage[mixed]{preamble}
\usepackage{lecture}
% \lecturelevel{chapter}  % book sınıfında varsayılan: \topic -> \chapter

\title{DGS Mathematics}
\author{math-lab}
\date{\today}

\begin{document}
\frontmatter
\maketitle

\chapter*{Preface}
Bu kitap Dikey Geçiş Sınavı (DGS) sayısal bölümünün matematik ve geometri
konularını ispat tarzında anlatır. Konu anlatımları \texttt{topics/}, sorular
\texttt{problems/} altında tutulur ve bu kitap yalnızca onları sıralar; aynı
konular TYT ve AYT kitaplarında da ortak kullanılır.

\tableofcontents

\mainmatter

% --- Yol haritası ------------------------------------------------------------
% Konu yazıldıkça ilgili satırın yorumunu kaldırın. Her konudan sonra o konunun
% soruları \input edilir, ör.
%   \input{problems/number-theory/elementary/divisibility/dgs-2024-07}
%
% Temel kavramlar ve sayılar
% \input{topics/number-theory/elementary/integers}          % Temel kavramlar, sayı kümeleri
% \input{topics/number-theory/elementary/digits}            % Sayı basamakları
% \input{topics/number-theory/elementary/divisibility}      % Bölme ve bölünebilme
% \input{topics/number-theory/elementary/gcd-lcm}           % EBOB–EKOK
% \input{topics/algebra/elementary/rational-numbers}        % Rasyonel ve ondalık sayılar
% \input{topics/algebra/elementary/inequalities}            % Basit eşitsizlikler
% \input{topics/algebra/elementary/absolute-value}          % Mutlak değer
% \input{topics/algebra/elementary/exponents}               % Üslü sayılar
% \input{topics/algebra/elementary/radicals}                % Köklü sayılar
% \input{topics/algebra/elementary/factoring}               % Çarpanlara ayırma
% \input{topics/algebra/elementary/ratio-proportion}        % Oran–orantı
% \input{topics/algebra/elementary/linear-equations}        % Denklem çözme
%
% Problemler
% \input{topics/algebra/elementary/word-problems}           % Sayı, kesir, yaş, işçi, hareket,
%                                                           % yüzde–kâr–zarar, karışım
%
% Kümeler, fonksiyonlar, işlemler
% \input{topics/foundations/sets/sets}                      % Kümeler
% \input{topics/foundations/sets/functions}                 % Fonksiyonlar
% \input{topics/algebra/elementary/operations}              % İşlem
% \input{topics/number-theory/elementary/modular-arithmetic} % Modüler aritmetik
%
% Sayma ve olasılık
% \input{topics/combinatorics/enumerative/counting}         % Permütasyon, kombinasyon
% \input{topics/probability-statistics/elementary/probability} % Olasılık
% \input{topics/probability-statistics/elementary/data}     % Tablo ve grafik yorumlama
%
% Geometri
% \input{topics/geometry/euclidean/angles}                  % Doğruda ve üçgende açılar
% \input{topics/geometry/euclidean/triangles}               % Özel üçgenler, benzerlik, alan
% \input{topics/geometry/euclidean/polygons}                % Çokgenler ve dörtgenler
% \input{topics/geometry/euclidean/circles}                 % Çember ve daire
% \input{topics/geometry/analytic/plane}                    % Analitik geometri
% \input{topics/geometry/solid/solids}                      % Katı cisimler

\backmatter
% Kaynakça ilk \cite ile birlikte açılır:
% \bibliographystyle{plain}
% \bibliography{refs}

\end{document}
